\scanpage{41}
\[
\begin{aligned}
 gug^{-1}&=u+\lambda\\
 \operatorname{Tr}u&=\operatorname{Tr}u+2\lambda
\end{aligned}
\]
si 2 inv. ceci fait $\lambda=0$ on gagne.

Sinon \deleted{écrivons aussi} \deleted{il y a contreexemple,}
\[
\begin{aligned}
 \deleted{\det u=\det(u+\lambda)}\\
 \deleted{\phantom{\det u={}}\lambda^2+(\operatorname{Tr}u)\lambda+\det u}\\
 \deleted{\text{donc }\lambda\unsure{{}^{1}|{}|{}^{2}}(\operatorname{Tr}u+\lambda)=0}
\end{aligned}
\]
\deleted{par ex. la matrice diagonale \unsure{$\xi'$} $u=(\zeta,\zeta+\lambda)$ est équivalente à la matrice (diag) \deleted{\unsure{$\zeta,\zeta+1$}} $u+\lambda\mathrm{id}=(\zeta+\lambda,\zeta)$, \unsure{?} sans lui être égale ! Ici le Centralisateur de $\mathfrak h$ \deleted{dans} $G$ \deleted{\unsure{en}} (en car \unsure{2|$p$}) est lisse de \deleted{dim} relative 1 ! \deleted{\unsure{lisse de}}}

Si $u$ est diagonalisable, soit $u=(\zeta,\zeta')$ alors la translation par $\zeta+\zeta'$ le transforme en $(\zeta',\zeta)$, qui lui est équivalente.

\margin{?}Mais le centralisateur \inserted{\uline{connexe}} semble être \unsure{encore|tjs|donc} $H$.

\deleted{Non} \deleted{\unsure{car si on}} \uline{NB} Dans l'algèbre de Lie de $G$, \inserted{avec la base habituelle} \deleted{le centralisateur de}
\[
 [H,X]=X,\quad [H,Y]=-Y,\quad [X,Y]=2H\underset{\text{(en car 2}}{=0}
\]
le centralisateur de $X$\deleted{$\alpha H+\beta$} \deleted{\unsure{$xX+hH+yY$}} \deleted{est formé des $X$} est formé des $xX+hH+yY$ tels que $h=0$, donc est de dim trop grande -- celui de $H$ est o.k.\ldots
% archive sheet number: 40
